\scanpage{11}
% archive foliation: 10
\diagram{\begin{tikzcd}X\arrow[d,"f"']\\S\end{tikzcd}}
propre lisse dim rel. 1 fibres connexes
\[
\Omega^1_{X/S}=\omega\qquad f_*(\omega)=\Omega\qquad
\mathbf{R}^1f_*(\Omega^*_{X/S})=\mathcal{DR}
\]
\[
\mathrm{R}^1f_*(\mathcal{O}_X)\simeq\check\Omega\simeq t
\]
\diagram{\begin{tikzcd}[column sep=small]
0\arrow[r]&\Omega\arrow[r]&\mathcal{DR}\arrow[r]&\check\Omega\arrow[r]&0
\end{tikzcd}}
\margin{autodualité
\deleted{$\mathcal{DR}\times\mathcal{DR}\longrightarrow\mathcal{O}_S$ \uline{alternée}}}
\deleted{le faisceau des \unsure{splittings} \unsure{de} $\mathcal{DR}\to\check\Omega\to0$ \unsure{est un espace}}
\diagram{\begin{tikzcd}[column sep=small]
\deleted{0}\arrow[r]&\deleted{\unsure{\check\Omega\otimes\check\Omega}}\arrow[r]&\deleted{\unsure{\Sigma}}\arrow[r]&\deleted{\mathcal{O}_S}\arrow[r,dash]&\deleted{0}
\end{tikzcd}}
\[
\mathcal{DR}\times\mathcal{DR}\xrightarrow[\varphi]{\text{alterné}}\mathcal{O}_S
\qquad(=\mathbf{R}^2f_*(\Omega^*_{X/S})\simeq\mathrm{R}^1f_*(\Omega^1_{X/S}))
\]
\deleted{nulle sur $\Omega\times\Omega$, d'où}
l'orthogonal de $\Omega$ est $\Omega$, :
\[
\deleted{\unsure{\mathcal{DR}}\quad\Omega\otimes\check\Omega}
\qquad\bigwedge^2\mathcal{DR}\xrightarrow[\varphi]{}\mathcal{O}_S
\]
\diagram{\begin{tikzcd}[column sep=small]
0\arrow[r,dash]&\Omega\arrow[r,dash]\arrow[d,"\mathrm{id}"']&\mathcal{DR}\arrow[r,dash]\arrow[d,"\unsure{\sim}"]&\check\Omega\arrow[r,dash]\arrow[d,"\mathrm{id}"]&0\\
0\arrow[r,dash]&\Omega\arrow[r,dash]&\check{\mathcal{DR}}\arrow[r,dash]&\check\Omega\arrow[r]&0
\end{tikzcd}}
\unsure{I.e.} on connaît déjà l'accouplement sur $\Omega\times\mathcal{DR}$ et $\mathcal{DR}\times\Omega$ --- l'indétermination est dans \deleted{\unsure{$\Omega$}} $\operatorname{Hom}(\check\Omega\otimes\check\Omega,\mathcal{O}_S)=\Omega\otimes\Omega$ et $=$ (puisqu'il s'agit de formes alternées) dans $\bigwedge^2\Omega$
\diagram{\begin{tikzcd}[column sep=large]
0\arrow[r,dash,"\bigwedge^2\Omega"]&\uline{\mathrm{Fil}^2}\arrow[r,dash,"\Omega\otimes\check\Omega"]&\mathrm{Fil}^1\arrow[r,dash,"\bigwedge^2\check\Omega"]\arrow[dr,"\unsure{\text{éval}}"']&\bigwedge^2\mathcal{DR}\arrow[d,dashed,"\varphi"]\\
&&&\mathcal{O}_S
\end{tikzcd}}
Pour un splitting de $\mathcal{DR}$ par $q:\check\Omega\to\mathcal{DR}$, on aura comme matrice de $\varphi$ pour la décomposition $\mathcal{DR}\overset{q}{\simeq}\Omega\oplus\check\Omega$
\diagram{\[
\begin{pmatrix}0&+\mathrm{id}\\-\mathrm{id}&\varphi_0\end{pmatrix}
\]}
où $\varphi_0$ est une forme alternée sur $\check\Omega$. Pour un autre choix $q'(t)=q(t)+u(t)$ ($u:\check\Omega\to\Omega$) on aura
\[
\begin{aligned}
\varphi'_0(t,t')&=\varphi(q(t)+u(t),q(t')+u(t'))\\
&=\varphi_0(t,\unsure{t'})\unsure{-|+}\langle t,u(t')\rangle\unsure{+|-}\langle t',u(t)\rangle\deleted{\unsure{+}}
\end{aligned}
\]
